\section{The Scaffold}\label{sec:scaffold}

A scaffold, as the introduction said, is code that people write and
check once, and that fixes the form of every theorem built on it. Ours
has two parts: the library of our earlier paper, which fixes the
syntactic side, and the additions of this work, which fix the semantic
side. This section describes both, says which recurring difficulties
the additions remove, and ends with what a theorem looks like as a
result. Table~\ref{tab:scaffold} lists the components.

\subsection{What the library fixes}\label{sec:library}

The library~\cite{anonymous2027popl} represents a relation set as an
inductive family \aid{Γ : WRel X} of pairs of words, and the congruence
it generates as an inductive family \aid{\_≈\_} with one constructor per
closure rule and one for the axioms. A presentation theorem is a record.
\aid{Γ IsSubPresentationOf G} holds an interpretation \aid{⟦\_⟧} of
words in the group \aid{G}, with a proof that it is a group
homomorphism from the words modulo \aid{≈} and that it is injective,
that is, sound and complete; \aid{Γ IsPresentationOf G} adds that it is
onto. Both are stated for any group, or monoid, of the standard
library's algebra hierarchy. Around them the library has morphisms of
presentations, normal forms as the standard library's function bundles
with the lemma that soundness and a unique normal form give
completeness, a Reidemeister--Schreier engine of coset tables and
towers, and presentations of direct, semidirect, $n$-fold and
amalgamated products and of group extensions given a short exact
sequence, each with a theorem that the constructed relation set
presents the constructed group.

For circuits it fixes the syntax. A gate set is a family
\aid{Gate : ℕ → Set} indexed by arity; a generator is a gate on the
bottom wires, or a generator shifted up one wire; a circuit on $n$
wires is a word over the generators. \aid{Lift-Relation} extended any
gate relation with three structural rules, \aid{cong↑}, \aid{comm₁} and
\aid{comm₂}; this work added the gates on no wires, the scalars, with
their rule \aid{ω↑=ω} (Section~\ref{sec:presentations}), so that a
paper's figure of rules, each stated at the width it is drawn on, is
the whole relation set.

Its semantics was abstract, and that is where its examples stopped. The
symmetric groups were presented as permutations, the wreath products
as signed permutations and the Pauli groups as $(\Z_p\times\Z_p)^n$;
the single-qubit and single-qutrit Clifford+$T$ monoids and
$U_3(\Z[\frac12,i])$ ended in an isomorphism between the gate-set
presentation and an amalgamated product of smaller presentations, with
the matrix monoid taken from the literature. Its discussion of
limitations said why: the matrix side needs rings and unitary groups
that the standard library does not have, and then the completeness
proof of each paper has to be formalized. Both are done here.

\subsection{What this work added}\label{sec:extensions}

The library proper, the folders \texttt{Word}, \texttt{Presentation},
\texttt{Normalization}, \texttt{Circuit} and \texttt{ForStdlib}, is
now about 16{,}000 lines, against about 11{,}000 at the earlier paper.
The additions are of four kinds; the first sits in the developments'
folders and is counted with them.

\paragraph{The matrix layer.}\label{sec:matrix-layer} The layer has two
ends. At the matrix end are square matrices indexed by $\mathsf{Fin}\,n$
over the rings of exact synthesis, $\D[\omega]=\Z[\frac1{\sqrt2},i]$,
$\D[i]=\Z[\frac12,i]$, $\D[\sqrt2]=\Z[\frac1{\sqrt2}]$ and
$\D=\Z[\frac12]$, taken from EucDomain, an Agda port of a Haskell
library for exact synthesis. Over each ring the unitary group $U_n$, or
the orthogonal group $O_n$ when the involution is the identity, is a
group of the standard library; the one- and two-level generators act on
columns; and embeddings between dimensions move a relation on $d$
indices to any $d$ indices of a larger matrix. The ring operations of
$\D[\omega]$, $\D[\sqrt2]$ and $\D$ are opaque, and over $\D[i]$ the
vector updates are, so that no proof unfolds the arithmetic of a
dyadic fraction unless it asks to (Section~\ref{sec:engineering}). At
the circuit end are operators on bit vectors, functions of two vectors
of $n$ bits over any commutative ring with an involution, with sums
over bit vectors, the Kronecker product and its mixed product law, a
gate placed on the bottom wires of a wider circuit, and an operator
moved up a wire. This is what makes the structural rules sound once
for every gate set: a gate on the bottom wires commutes with a shifted
operator by the mixed product law, whatever the gate. A map from
operators to matrices, which renumbers a sum over bit vectors as a sum
over $\mathsf{Fin}\,2^n$, is a monoid homomorphism, and it sends an
operator whose conjugate transpose is its inverse to an element of
$U_{2^n}$. The circuit end exists in three versions: over an abstract
ring, as just described; over $\Z[\sqrt2]$ with a power of
$\frac1{\sqrt2}$ kept apart, with stored tries and a decision procedure
on row operations, for the developments whose soundness checks are
large (Section~\ref{sec:eval}); and over the cyclotomic rings of
qudits, with a comparison up to a phase, for the minimality theorems.
The first two reach the groups at the matrix end; the qudit version
stops at scaled operators compared up to a phase, which is all the
minimality theorems need. Each circuit development chooses a ring and
gives its gate matrices, which with the relations is all a reader has
to compare with the paper. The layer is about 6{,}000 lines, counted in
the developments' totals of Table~\ref{tab:corpus} rather than in the
library's. It grew with the developments, and its modules live in the
folders of the developments that first needed them: at the matrix end
\texttt{Ring}, \texttt{Vector}, \texttt{Algebra},
\texttt{MatrixAlgebra}, \texttt{Unitary}, \texttt{Action},
\texttt{Embedding} and \texttt{EmbedAction} of
\texttt{Clifford+CS-TwoLevel}, with the \texttt{Semantics} module of
each two-level development; at the circuit end the \texttt{Semantics}
folders of \texttt{Real-Clifford}, \texttt{Real-Clifford+CH} and
\texttt{Examples/Independence}. They are imported across developments;
we have not yet moved them into one folder.

\paragraph{Independence.}\label{sec:independence-layer} A rule of a
presentation is independent when the congruence generated by the other
rules does not relate its two sides. The classical proof is a model in
which the other rules hold and this one fails, and
\aid{Presentation.Independence} proves that argument once: a selection
of axiom instances, a reading of words in a monoid that respects the
selected instances, and the conclusion that the reading respects
everything they derive, so that an instance it separates is
independent. \aid{Circuit.Independence} does the same for circuits,
where a model is a monoid per width related by a shift homomorphism,
and uses it on the structural rules: each of the four is independent
of the other three, and a fifth rule that an earlier version of this
work had, that a scalar commutes with every gate, is derivable from
the four when scalars commute with one another, and not otherwise. On
top of these, the modules for Blake's minimality
theorems~\cite{blake2026simpler} add the swap as a generator with the
rules of a PROP, the notion of a fragment with named rules, the fact
that a derivation at width $n$ uses no rule of arity above $n$ (his
Remark~14), and the weighted models of his Figure~8, so that a
minimality theorem is one model per rule, each checked at its arity by
\aid{refl}.

\paragraph{Routes, tactics and translations.} Two of the proof routes
of Section~\ref{sec:architecture} are module shapes written once: the
levels and edges of exact synthesis, with completeness by well-founded
induction from a Main Lemma, and paths stated over any normal-word
function; and the back-and-forth argument of an encoding. So is the
rewriter that checks a derivation given as a list of steps. Each
development has its own copy of these, with the parameters changed.
The other two routes, normal forms and induction on the width, are
patterns that the author follows with the library's records, and the
product constructions are the library's. Two additions are library
modules proper. Bian and Selinger's tactic
library~\cite{bian2023cliffordt2,bian2023cliffordcs3}, about 2{,}800
lines, is rebased on the library's words and congruences, so that their
computer-generated proofs port mechanically and the rewriters of four
developments, one of them shared by a fifth, can use its canonical
forms for commuting letters. And \aid{Presentation.Translation} turns a
sub-presentation of a group and a presentation of the same group into a
translation of words that preserves the matrix, which is how a real
Clifford circuit becomes a two-level word
(Section~\ref{sec:theorem-form}).

\paragraph{Constructions, and the standard library.} The direct,
semidirect and amalgamated products and the extensions by a short
exact sequence are the earlier library's. Extensions by a factor set
built from data on the generators were added for the qudit Clifford
development, together with what it needed of the standard library:
actions of monoids and groups, $\mathbb{N}$-indexed families of groups
with embeddings between levels, central extensions presented by a
cocycle, and $\Z/n\Z$ as $\mathsf{Fin}\,n$ with its ring, moved from
the earlier library's own folder and given a ring solver. Central
products were added as well, and no development uses them yet. These
are written in the standard library's style, as the earlier additions
were.

\paragraph{What the additions remove.} Five things recur in these
papers, and each is done once here rather than once per paper: the
soundness of the structural rules, a proof in the Kronecker calculus;
the soundness of the relations, a decision procedure, with stored
tries where the check is large; the independence of a rule, a model and
one lemma; the induction on levels, a module with completeness from a
Main Lemma by one lemma that closes a square, and the encoding
argument, a module of a hundred lines; and the identity of a group
with the group of another paper, which here is literal: the real
stabilizer circuits, the real Clifford+$CH$ circuits and the two-level
presentations of Hadamard-$\Pi$ and of Cl\'ement are stated against the
same $O_{2^n}(\Z[\frac1{\sqrt2}])$, and the Clifford, CNOT-dihedral and
real stabilizer circuits against the same $U_{2^n}(\D[\omega])$, so
that a theorem relating two developments can be stated, and where one
is proved (Section~\ref{sec:encoding}) the developments check each
other.

\begin{table}
\caption{The scaffold. ``Earlier'' is the library of
\citet{anonymous2027popl}, ``here'' this work. The matrix layer's
modules sit in the folders of the developments that first needed them.}
\label{tab:scaffold}
\footnotesize
\setlength{\tabcolsep}{4pt}
\begin{tabularx}{\linewidth}{@{}Y l Y Y@{}}
\toprule
Component & Origin & Where & Used by \\
\midrule
Words, congruences, the presentation records, morphisms & earlier & \texttt{Word}, \texttt{Presentation} & all \\
Circuits by width with three structural rules; scalars with their rule, and a scalar's centrality as a theorem & earlier; here & \texttt{Circuit} & circuit developments \\
Normal forms as bundles; completeness from soundness and uniqueness & earlier & \texttt{Normalization} & qudit Clifford, real Clifford+$CH$ \\
Reidemeister--Schreier engine, coset tables and towers & earlier & \texttt{Normalization} & qudit Clifford, real Clifford+$CH$, the ports \\
Products and extensions of presentations; extensions by a factor set & earlier; here & \texttt{Presentation/\allowbreak Construct} & qudit Clifford \\
Bian--Selinger's tactics, with their Reidemeister--Schreier interface & here & \texttt{Presentation/\allowbreak Tactics} & the ports; four rewriters \\
Translation along a presentation of the same group & here & \texttt{Presentation/\allowbreak Translation} & real stabilizer into Hadamard-$\Pi$ \\
Independence by models; the structural rules; fragments in a PROP, weighted models & here & \texttt{Presentation}, \texttt{Circuit}, \texttt{Examples/\allowbreak Independence} & minimality; the structural rules \\
Matrices over $\mathsf{Fin}\,n$ and EucDomain's rings; $U_n$, $O_n$; two-level generators; embeddings & here & \texttt{Clifford+CS-\allowbreak TwoLevel}; each two-level \texttt{Semantics} & two-level developments; the group of every circuit theorem \\
Operators on bit vectors; Kronecker calculus; the map to matrices; stored tries, decision procedure; qudits up to a phase & here & \texttt{Semantics} of \texttt{Real-Clifford}, \texttt{Real-Clifford+CH}, \texttt{Independence} & real stabilizer, Clifford, real Clifford+$CH$, minimality \\
Route shapes: levels and edges, paths over any normal word, back-and-forth encoding, rewriter & here & one copy per development & two-level, real Clifford+$CH$, five with derivations \\
Actions, indexed families of groups, central extensions, factor sets, $\Z/n\Z$ with a solver & here & \texttt{ForStdlib} & qudit Clifford, phase-affine \\
\bottomrule
\end{tabularx}
\end{table}

\subsection{The form of a theorem}\label{sec:theorem-form}

Every development ends in a module that states its theorem in the
scaffold's terms, without hypotheses; the one exception is the
Hadamard-$\Pi$ folder, whose theorem takes its hard case as a
hypothesis that the Cl\'ement folder discharges, so the unconditional
statement lives there. The real stabilizer development concludes
\begin{agdacode}
presentation : {n : ℕ} → (n VRel,_===_) IsSubPresentationOf U (2 ^ n)
\end{agdacode}
The statement has two parameters. The relation \aid{n VRel,\_===\_} is
the paper's R1--R16 lifted to $n$ wires by the structural rules, and
the type of the record requires it to be a relation on the library's
circuits. The group \aid{U (2 \textasciicircum{} n)} is the unitary
group of the layer, here over any nontrivial commutative ring $A$ with
an involution $x\mapsto\overline{x}$, playing the role of complex
conjugation, and an element $s$ with $2s^2=1$ and $\overline{s}=s$,
playing the role of $\frac1{\sqrt2}$; for example $\mathbb{C}$, or
$\Z[\frac1{\sqrt2}]$. A reader who has seen one such statement reads
the others the same way: \aid{(\_===\_ \{n\}) IsPresentationOf (O n)}
for the two-level groups, \aid{Rel IsSubPresentationOf U 8} for the
three-qubit Clifford+$CS$ circuits inside $U_8(\D[i])$, and so on down
Table~\ref{tab:corpus}.

To trust the line above, a reader checks four things, and nothing else
in the 11{,}799 lines of the development. The relation: the module that
states it has 117 lines, and Section~\ref{sec:example} shows part of
it. The group: it is built from the operators and the unitary matrices
of the layer, whose definitions take under 200 lines and whose lemmas
need not be read, and a reader who has checked the layer once need not
check it again for the next paper. The interpretation: it sends $-1$,
$H$, $Z$ and $CZ$ to their matrices and a word to the product of its
letters, in about 30 lines. And the one line of the theorem module that
fills the record's field \aid{⟦\_⟧} with that interpretation: the record
holds the interpretation as a field, so \aid{R IsSubPresentationOf U}
on its own says that the presented group embeds in \aid{U} in some way,
and the reader looks at the field to see in which. These come to under
500 lines, most of them the layer's definitions. The rest of the
library, which defines words, circuits, the structural rules and the
record itself, is checked once for all the developments.

A second module instantiates this theorem at
$\D[\sqrt2]=\Z[\frac1{\sqrt2}]$ and $s=\frac1{\sqrt2}$, which puts the
real stabilizer group inside the orthogonal group of
Hadamard-$\Pi$. Because that group is presented by Figure~6, every real
Clifford circuit translates to a two-level word with the same matrix,
and two circuits are equal by the circuit relations exactly when their
translations are equal by Figure~6. The two developments, written
separately, thus check each other. A third instantiates it at
$\D[\omega]$, into the unitary group that the Clifford and CNOT-dihedral
developments use.
